\section{Mô hình hóa Mối kết hợp \& Bản số Dữ liệu}

\begingroup
\setlength{\tabcolsep}{3.5pt} % Thu nhỏ khoảng đệm giữa các cột
\small % Cỡ chữ vừa vặn giúp bảng đẹp và cân đối

\begin{longtable}{|
>{\raggedright\arraybackslash}p{0.24\textwidth}|
>{\raggedright\arraybackslash}p{0.25\textwidth}|
>{\raggedright\arraybackslash}p{0.17\textwidth}|
>{\raggedright\arraybackslash}p{0.34\textwidth}|}
\hline

\textbf{Mối quan hệ \& Phân loại thực thể} &
\textbf{Bản số \& Mức tham gia} &
\textbf{Thuộc tính mối quan hệ} &
\textbf{Mô tả nghiệp vụ} \\
\hline
\endfirsthead

\hline
\textbf{Mối quan hệ \& Phân loại thực thể} &
\textbf{Bản số \& Mức tham gia} &
\textbf{Thuộc tính mối quan hệ} &
\textbf{Mô tả nghiệp vụ} \\
\hline
\endhead

\textbf{User — Workspace — Role} (\texttt{<participates>}) \newline
\textit{Mối kết hợp 3 ngôi:} User (Mạnh), Workspace (Mạnh), Role (Mạnh) &
\textbf{Tỷ lệ:} M:N:1 \newline
\textbf{Bản số:} User (0..*), Workspace (1..*), Role (0..*) \newline
\textbf{Tham gia:} User (Partial 0), Workspace (Total 1), Role (Partial 0) &
\texttt{joined\allowbreak\_at}, \texttt{status} &
Ghi nhận sự tham gia của Người dùng vào Workspace. Trong từng Workspace, mỗi user bắt buộc gán đúng một Role (Scoped RBAC). Một user có thể thuộc nhiều Workspace với vai trò độc lập. \\ \hline

\textbf{Role — Permission} (\texttt{<contains>}) \newline
Role (Mạnh) $\leftrightarrow$ Permission (Mạnh) &
\textbf{Tỷ lệ:} M:N \newline
\textbf{Bản số:} Role (0..*) $\leftrightarrow$ Permission (0..*) \newline
\textbf{Tham gia:} Role (Partial) $\leftrightarrow$ Permission (Partial) &
\texttt{granted\allowbreak\_at} &
Một vai trò chứa tập hợp nhiều quyền; một quyền hạn có thể thuộc về nhiều vai trò khác nhau. \\ \hline

\textbf{User — Github\allowbreak\_Connection} (\texttt{<connects>}) \newline
User (Mạnh) $\leftrightarrow$ Github\_Connection (Phụ thuộc User) &
\textbf{Tỷ lệ:} 1:N \newline
\textbf{Bản số:} User (0..*), Github\allowbreak\_Connection (1..1) \newline
\textbf{Tham gia:} User (Partial), Github\allowbreak\_Connection (Total) &
Không có &
Ghi nhận việc User ủy quyền tài khoản GitHub qua OAuth2, lưu \texttt{access\_token} và scopes để đồng bộ repo, commit, PR, webhook. \\ \hline

\textbf{Workspace — Repository} (\texttt{<owns>}) \newline
Workspace (Mạnh) $\leftrightarrow$ Repository (Mạnh) &
\textbf{Tỷ lệ:} 1:N \newline
\textbf{Bản số:} Workspace (1..1) $\leftrightarrow$ Repository (0..*) \newline
\textbf{Tham gia:} Workspace (Partial), Repository (Total) &
Không có &
Một Workspace quản lý nhiều Repository; mỗi Repository thuộc quyền sở hữu của một Workspace. \\ \hline

\textbf{Repository — Branch} (\texttt{<has\_branch>}) \newline
Repository (Mạnh) $\leftrightarrow$ Branch (Yếu - Phụ thuộc Repo) &
\textbf{Tỷ lệ:} 1:N \newline
\textbf{Bản số:} Repository (1..1) $\leftrightarrow$ Branch (1..*) \newline
\textbf{Tham gia:} Repository (Total), Branch (Total) &
Không có &
Một kho chứa mã nguồn có ít nhất một nhánh mặc định; mỗi nhánh bắt buộc nằm trong duy nhất một kho chứa. \\ \hline

\textbf{Branch — Commit} (\texttt{<contains\_commit>}) \newline
Branch (Yếu) $\leftrightarrow$ Commit (Yếu - Phụ thuộc Branch/Repo) &
\textbf{Tỷ lệ:} 1:N \newline
\textbf{Bản số:} Branch (0..*) $\leftrightarrow$ Commit (1..1) \newline
\textbf{Tham gia:} Branch (Partial), Commit (Total) &
Không có &
Một nhánh chứa nhiều lượt commit; mỗi commit thuộc về đúng một nhánh phát triển. \\ \hline

\textbf{User — Commit} (\texttt{<makes\_commit>}) \newline
User (Mạnh) $\leftrightarrow$ Commit (Yếu) &
\textbf{Tỷ lệ:} 1:N \newline
\textbf{Bản số:} User (0..*) $\leftrightarrow$ Commit (1..1) \newline
\textbf{Tham gia:} User (Partial), Commit (Total) &
Không có &
Một người dùng thực hiện nhiều commit; mỗi commit được tạo bởi đúng một người dùng. \\ \hline

\textbf{Repository — Pull\allowbreak Request} (\texttt{<receives\_pr>}) \newline
Repository (Mạnh) $\leftrightarrow$ PullRequest (Yếu - Phụ thuộc Repo) &
\textbf{Tỷ lệ:} 1:N \newline
\textbf{Bản số:} Repository (0..*) $\leftrightarrow$ PullRequest (1..1) \newline
\textbf{Tham gia:} Repository (Partial), PullRequest (Total) &
Không có &
Một Repository chứa nhiều Yêu cầu hợp nhất (PR); mỗi PR gửi về đúng một kho chứa. \\ \hline

\textbf{Commit — Code\allowbreak Diff} (\texttt{<generates\_diff>}) \newline
Commit (Yếu) $\leftrightarrow$ CodeDiff (Yếu - Phụ thuộc Commit/PR) &
\textbf{Tỷ lệ:} 1:N \newline
\textbf{Bản số:} Commit (1..*) $\leftrightarrow$ CodeDiff (1..1) \newline
\textbf{Tham gia:} Commit (Total), CodeDiff (Total) &
Không có &
Một Commit sinh ra một hoặc nhiều bản ghi khác biệt code (Diff); mỗi bản ghi diff thuộc về duy nhất một commit. \\ \hline

\textbf{Pull\allowbreak Request — Commit} (\texttt{<includes\_commit>}) \newline
PullRequest (Yếu) $\leftrightarrow$ Commit (Yếu) &
\textbf{Tỷ lệ:} 1:N \newline
\textbf{Bản số:} PullRequest (1..*) $\leftrightarrow$ Commit (1..1) \newline
\textbf{Tham gia:} PullRequest (Total), Commit (Partial) &
Không có &
Một Yêu cầu hợp nhất (PR) tổng hợp danh sách gồm một hoặc nhiều commit; mỗi commit thuộc về đúng một PR. \\ \hline

\textbf{Commit — File} (\texttt{<modifies>}) \newline
Commit (Yếu) $\leftrightarrow$ File (Yếu - Phụ thuộc Repo) &
\textbf{Tỷ lệ:} M:N \newline
\textbf{Bản số:} Commit (1..*) $\leftrightarrow$ File (0..*) \newline
\textbf{Tham gia:} Commit (Total), File (Partial) &
\texttt{change\allowbreak\_type}, \newline
\texttt{lines\allowbreak\_added}, \newline
\texttt{lines\allowbreak\_deleted} &
Một Commit tác động chỉnh sửa ít nhất một File; một File có thể chịu sự biến đổi qua nhiều lượt Commit. \\ \hline

\textbf{Repository — File} (\texttt{<contains\_file>}) \newline
Repository (Mạnh) $\leftrightarrow$ File (Yếu - Phụ thuộc Repo) &
\textbf{Tỷ lệ:} 1:N \newline
\textbf{Bản số:} Repository (0..*) $\leftrightarrow$ File (1..1) \newline
\textbf{Tham gia:} Repository (Partial), File (Total) &
Không có &
Một Repository quản lý cấu trúc cây thư mục chứa nhiều File mã nguồn; mỗi File thuộc một kho chứa duy nhất. \\ \hline

\textbf{File — Code\allowbreak Entity} (\texttt{<parses\_into>}) \newline
File (Yếu) $\leftrightarrow$ CodeEntity (Yếu - Phụ thuộc File) &
\textbf{Tỷ lệ:} 1:N \newline
\textbf{Bản số:} File (0..*) $\leftrightarrow$ Code\allowbreak Entity (1..1) \newline
\textbf{Tham gia:} File (Partial), Code\allowbreak Entity (Total) &
Không có &
Một File code qua bóc tách AST thu được nhiều thực thể code (Class, Function); mỗi thực thể code nằm trọn trong đúng một File. \\ \hline

\textbf{Code\allowbreak Entity — Code\allowbreak Entity} (\texttt{<depends\_on>}) \newline
\textit{Đệ quy:} CodeEntity (Yếu) $\leftrightarrow$ CodeEntity (Yếu) &
\textbf{Tỷ lệ:} M:N \newline
\textbf{Bản số:} Phụ thuộc (0..*), Được phụ thuộc (0..*) \newline
\textbf{Tham gia:} Cả 2 phía Partial &
\texttt{dependency\allowbreak\_type} &
Xác định mối quan hệ phụ thuộc (gọi hàm, kế thừa, import) giữa các đơn vị code trong hệ thống. \\ \hline

\textbf{Workspace — Template} (\texttt{<defines>}) \newline
Workspace (Mạnh) $\leftrightarrow$ Template (Mạnh) &
\textbf{Tỷ lệ:} 1:N \newline
\textbf{Bản số:} Workspace (0..*) $\leftrightarrow$ Template (1..1) \newline
\textbf{Tham gia:} Workspace (Partial), Template (Total) &
Không có &
Một Workspace định nghĩa nhiều khuôn mẫu tài liệu; mỗi mẫu tài liệu thuộc sở hữu của một Workspace. \\ \hline

\textbf{Template — Template\allowbreak Section} (\texttt{<has\_section>}) \newline
Template (Mạnh) $\leftrightarrow$ TemplateSection (Yếu - Phụ thuộc Template) &
\textbf{Tỷ lệ:} 1:N \newline
\textbf{Bản số:} Template (1..*) $\leftrightarrow$ Template\allowbreak Section (1..1) \newline
\textbf{Tham gia:} Template (Total), Template\allowbreak Section (Total) &
Không có &
Một Template cấu thành từ ít nhất một Section dàn ý; mỗi Section nằm trên đúng một Template. \\ \hline

\textbf{Template\allowbreak Section — Placeholder} (\texttt{<contains\_placeholder>}) \newline
TemplateSection (Yếu) $\leftrightarrow$ Placeholder (Yếu - Phụ thuộc Section) &
\textbf{Tỷ lệ:} 1:N \newline
\textbf{Bản số:} Template\allowbreak Section (0..*) $\leftrightarrow$ Placeholder (1..1) \newline
\textbf{Tham gia:} Template\allowbreak Section (Partial), Placeholder (Total) &
Không có &
Mỗi mục trong mẫu tài liệu chứa nhiều biến chèn tự động. Khi xóa Section, các Placeholder nằm trong đó bị xóa theo. \\ \hline

\textbf{Placeholder — Code\allowbreak Entity} (\texttt{<maps\_to>}) \newline
Placeholder (Yếu) $\leftrightarrow$ CodeEntity (Yếu) &
\textbf{Tỷ lệ:} M:N \newline
\textbf{Bản số:} Placeholder (0..*) $\leftrightarrow$ Code\allowbreak Entity (0..*) \newline
\textbf{Tham gia:} Placeholder (Partial), Code\allowbreak Entity (Partial) &
\texttt{extraction\allowbreak\_rule}, \newline
\texttt{is\allowbreak\_active} &
Ánh xạ chỉ thị cho AI biết lấy thông tin từ thành phần code nào để điền vào biến trong mẫu. \\ \hline

\textbf{Workspace — Document} (\texttt{<manages>}) \newline
Workspace (Mạnh) $\leftrightarrow$ Document (Mạnh) &
\textbf{Tỷ lệ:} 1:N \newline
\textbf{Bản số:} Workspace (0..*) $\leftrightarrow$ Document (1..1) \newline
\textbf{Tham gia:} Workspace (Partial), Document (Total) &
Không có &
Workspace quản lý danh mục tài liệu gốc của dự án; một tài liệu thuộc về đúng một Workspace. \\ \hline

\textbf{User — Document} (\texttt{<authors>}) \newline
User (Mạnh) $\leftrightarrow$ Document (Mạnh) &
\textbf{Tỷ lệ:} 1:N \newline
\textbf{Bản số:} User (0..*) $\leftrightarrow$ Document (1..1) \newline
\textbf{Tham gia:} User (Partial), Document (Total) &
Không có &
Người dùng đứng tên tác giả khởi tạo tài liệu; mỗi tài liệu có một tác giả duy nhất. \\ \hline

\textbf{Document — Document\allowbreak Version} (\texttt{<has\_version>}) \newline
Document (Mạnh) $\leftrightarrow$ DocumentVersion (Yếu - Phụ thuộc Document) &
\textbf{Tỷ lệ:} 1:N \newline
\textbf{Bản số:} Document (1..*) $\leftrightarrow$ Document\allowbreak Version (1..1) \newline
\textbf{Tham gia:} Document (Total), Document\allowbreak Version (Total) &
Không có &
Một tài liệu chính thức chứa ít nhất một phiên bản nội dung; mỗi phiên bản đại diện cho một mốc chỉnh sửa. \\ \hline

\textbf{Document\allowbreak Version — Change\allowbreak Log} (\texttt{<tracks\_change>}) \newline
DocumentVersion (Yếu) $\leftrightarrow$ ChangeLog (Yếu - Phụ thuộc DocVersion) &
\textbf{Tỷ lệ:} 1:N \newline
\textbf{Bản số:} Document\allowbreak Version (0..*) $\leftrightarrow$ ChangeLog (1..1) \newline
\textbf{Tham gia:} Document\allowbreak Version (Partial), ChangeLog (Total) &
Không có &
Mỗi phiên bản tài liệu có thể ghi nhận nhiều lượt nhật ký biến động nội dung chi tiết. \\ \hline

\textbf{Document\allowbreak Version — Review} (\texttt{<has\_review>}) \newline
DocumentVersion (Yếu) $\leftrightarrow$ Review (Yếu - Phụ thuộc DocVersion) &
\textbf{Tỷ lệ:} 1:N \newline
\textbf{Bản số:} Document\allowbreak Version (0..*) $\leftrightarrow$ Review (1..1) \newline
\textbf{Tham gia:} Document\allowbreak Version (Partial), Review (Total) &
Không có &
Một phiên bản tài liệu có thể nhận 0 hoặc nhiều lượt đánh giá kỹ thuật; mỗi lượt đánh giá thuộc về duy nhất một phiên bản. \\ \hline

\textbf{User — Review} (\texttt{<performs\_review>}) \newline
User (Mạnh) $\leftrightarrow$ Review (Yếu) &
\textbf{Tỷ lệ:} 1:N \newline
\textbf{Bản số:} User (0..*) $\leftrightarrow$ Review (1..1) \newline
\textbf{Tham gia:} User (Partial), Review (Total) &
Không có &
Một người dùng có thể thực hiện 0 hoặc nhiều lượt đánh giá; mỗi lượt review do đúng một người dùng chịu trách nhiệm. \\ \hline

\textbf{Review — Review\allowbreak Comment} (\texttt{<contains\_comment>}) \newline
Review (Sự kiện/Yếu) $\leftrightarrow$ ReviewComment (Yếu - Phụ thuộc Review) &
\textbf{Tỷ lệ:} 1:N \newline
\textbf{Bản số:} Review (0..*) $\leftrightarrow$ Review\allowbreak Comment (1..1) \newline
\textbf{Tham gia:} Review (Partial), Review\allowbreak Comment (Total) &
Không có &
Một phiên kiểm duyệt kỹ thuật chứa nhiều bình luận góp ý chi tiết; mỗi bình luận gắn chặt với một lượt review. \\ \hline

\textbf{Document\allowbreak Version — Approval} (\texttt{<requires\_approval>}) \newline
DocumentVersion (Yếu) $\leftrightarrow$ Approval (Yếu - Phụ thuộc DocVersion) &
\textbf{Tỷ lệ:} 1:N \newline
\textbf{Bản số:} Document\allowbreak Version (0..*) $\leftrightarrow$ Approval (1..1) \newline
\textbf{Tham gia:} Document\allowbreak Version (Partial), Approval (Total) &
Không có &
Một phiên bản tài liệu có thể có 0 hoặc nhiều lượt phê duyệt; mỗi lượt phê duyệt bắt buộc thuộc về đúng một phiên bản. \\ \hline

\textbf{User — Approval} (\texttt{<grants\_approval>}) \newline
User (Mạnh) $\leftrightarrow$ Approval (Yếu) &
\textbf{Tỷ lệ:} 1:N \newline
\textbf{Bản số:} User (0..*) $\leftrightarrow$ Approval (1..1) \newline
\textbf{Tham gia:} User (Partial), Approval (Total) &
Không có &
Một Manager có thể thực hiện 0 hoặc nhiều đợt phê duyệt; mỗi bản ghi phê duyệt do đúng một Manager ra quyết định. \\ \hline

\textbf{Document\allowbreak Version — Ground\allowbreak ing\allowbreak Evidence} (\texttt{<has\_evidence>}) \newline
DocumentVersion (Yếu) $\leftrightarrow$ GroundingEvidence (Yếu - Phụ thuộc DocVersion) &
\textbf{Tỷ lệ:} 1:N \newline
\textbf{Bản số:} Document\allowbreak Version (0..*) $\leftrightarrow$ Grounding\allowbreak Evidence (1..1) \newline
\textbf{Tham gia:} Document\allowbreak Version (Partial), Grounding\allowbreak Evidence (Total) &
Không có &
Phiên bản tài liệu lưu trữ bảng minh chứng đối soát nguồn gốc code để chống ảo giác (hallucination) của AI. \\ \hline

\textbf{Ground\allowbreak ing\allowbreak Evidence — Code\allowbreak Entity} (\texttt{<references\_code>}) \newline
GroundingEvidence (Yếu) $\leftrightarrow$ CodeEntity (Yếu) &
\textbf{Tỷ lệ:} M:N \newline
\textbf{Bản số:} Grounding\allowbreak Evidence (0..*) $\leftrightarrow$ Code\allowbreak Entity (0..*) \newline
\textbf{Tham gia:} Grounding\allowbreak Evidence (Partial), Code\allowbreak Entity (Partial) &
\texttt{similarity\allowbreak\_score}, \newline
\texttt{match\allowbreak\_type} &
Minh chứng AI liên kết trực tiếp với các đơn vị code cụ thể trong hệ thống để phục vụ đối soát sai lệch. \\ \hline

\textbf{Ground\allowbreak ing\allowbreak Evidence — Drift\allowbreak Alert} (\texttt{<detects>}) \newline
GroundingEvidence (Yếu) $\leftrightarrow$ DriftAlert (Thực thể sự kiện) &
\textbf{Tỷ lệ:} 1:N \newline
\textbf{Bản số:} Grounding\allowbreak Evidence (0..*) $\leftrightarrow$ DriftAlert (1..1) \newline
\textbf{Tham gia:} Grounding\allowbreak Evidence (Partial), DriftAlert (Total) &
Không có &
Minh chứng đối soát khi phát hiện sự bất đồng giữa tài liệu và mã nguồn mới sẽ kích hoạt cảnh báo sai lệch (Drift Alert). \\ \hline

\end{longtable}
\endgroup